\documentclass[12pt]{article}
\usepackage[T1]{fontenc}
\usepackage[utf8]{inputenc}
\usepackage{lmodern}
\usepackage{textcomp}
\usepackage{enumitem}
\usepackage{geometry}
\geometry{margin=1in}
\begin{document}
\begin{center}
Answer Key - Business Administration - Graduate Level Assessment 20 \
\small (Answers are ordered exactly as in the question paper.)
\end{center}

\section*{Multiple Choice}
\begin{enumerate}[label=\arabic*.]
\item \textbf{Answer:} A
\\ \textit{Options:} A) Price bundling.; B) Price differentiation.; C) Price skimming.; D) Price penetration.; E) Price discrimination.
\\ \textit{Note:} Price bundling occurs when competitors prioritize their pricing strategies based on the offerings of rivals rather than understanding and addressing customer needs, leading to a focus on competitive positioning rather than consumer value.
\item \textbf{Answer:} C
\\ \textit{Options:} A) \$850.28; B) \$846.50; C) \$848.72; D) \$845.00; E) \$850.00
\\ \textit{Note:} The correct answer is \$848.72 because the proceeds from the discounted draft are calculated by deducting the discount amount, which is based on the discount rate applied to the time remaining until maturity, from the face value of the draft.
\item \textbf{Answer:} D
\\ \textit{Options:} A) \$400.00; B) \$150.00; C) \$200.00; D) \$175.00; E) \$350.00
\\ \textit{Note:} Albert Morgan earned $175.00 because his commission is calculated by multiplying his sales of $3,500 by the 5\% commission rate, resulting in $3,500 x 0.05 = $175.00.
\item \textbf{Answer:} A
\\ \textit{Options:} A) \$6.00; B) \$1.50; C) \$5.00; D) \$4.00; E) \$0.20
\\ \textit{Note:} The millinerate is calculated by multiplying the cost per agate line by the circulation in millions, so $4 per line multiplied by 2 million equals $8 million, and when divided by 1 million, it results in the millinerate of \$6.00.
\item \textbf{Answer:} A
\\ \textit{Options:} A) 40\%, 0.25\%; B) 40\%, 1\%; C) 35\%, 0.75\%; D) 50\%, 1\%; E) 40\%, 0.5\%
\\ \textit{Note:} The correct answer is derived from converting the fractions to percentages by multiplying each by 100; (2/5) equals 40\% and (1/100) equals 1\%, thus (1/100) as a percentage is correctly represented as 0.25\%.
\end{enumerate}

\section*{True / False}
\begin{enumerate}[label=\arabic*.]
\setcounter{enumi}{5}
\item \textbf{Answer:} False
\\ \textit{Options:} A) True; B) False
\\ \textit{Note:} The statement is False because the total profit that can be distributed among partners cannot exceed the firm's net profit of $28,400, making it impossible for each partner to receive $10,800 if there are multiple partners.
\item \textbf{Answer:} False
\\ \textit{Options:} A) True; B) False
\\ \textit{Note:} The effective yield of a bond is calculated based on its coupon rate and purchase price, and for a bond with a coupon rate of 6.75\% purchased at 88.75, the actual yield would be higher than 8.2\%, making the statement false.
\item \textbf{Answer:} False
\\ \textit{Options:} A) True; B) False
\\ \textit{Note:} The statement is false because the probability of observing an earnings per share less than $5.5 would depend on the distribution of earnings per share and is not necessarily equal to 0.5000 unless the mean and median earnings per share are exactly $5.5 in a perfectly symmetrical distribution.
\item \textbf{Answer:} True
\\ \textit{Options:} A) True; B) False
\\ \textit{Note:} The statement is true because Robinson, Faris, and Wind (1967) indeed categorized organizational buying situations into three distinct types of buygroups, which serve as a foundational concept in understanding business-to-business purchasing behavior.
\item \textbf{Answer:} True
\\ \textit{Options:} A) True; B) False
\\ \textit{Note:} Replacing supply chains with supply loops highlights the importance of circular economy principles, which prioritize the recapture and reuse of waste materials in a firm's processes, ultimately promoting sustainability and resource efficiency.
\end{enumerate}

\section*{Long Form Response}
\begin{enumerate}[label=\arabic*.]
\setcounter{enumi}{10}
\item \textbf{Answer:} Mr. Glass's loan agreement with ABC Finance Company involves varying interest rates that significantly impact his total interest cost over the repayment period, which is calculated to be \$23. This figure arises from analyzing the terms of the loan, including the specific rates applied at different intervals and the principal amount borrowed. Such loan structures can complicate financial decision-making, as borrowers must carefully assess the total cost of borrowing rather than focusing solely on periodic payments. Additionally, understanding the implications of fluctuating interest rates is crucial for making informed financial choices, as it can affect cash flow and long-term financial planning. Ultimately, Mr. Glass's experience underscores the importance of thorough financial analysis in navigating complex loan agreements.
\\ \textit{Note:} correct because it accurately identifies the impact of varying interest rates on Mr. Glass's total interest cost, emphasizes the necessity of comprehensive analysis in loan agreements, and highlights the complexities these structures introduce to financial decision-making.
\item \textbf{Answer:} Karl Marx viewed the relationship between workers and bureaucracies as inherently exploitative, with bureaucracies serving the interests of capital rather than the workforce. He argued that bureaucratic structures prioritize efficiency and control over human needs, leading to alienation among workers who become mere cogs in a machine. This dynamic implies that contemporary business practices, which often emphasize hierarchical and rigid bureaucratic systems, may perpetuate worker discontent and disengagement. Therefore, organizations should consider adopting more flexible and participatory structures that empower employees, thereby aligning business goals with the well-being of the workforce and fostering a more engaged and productive environment.
\\ \textit{Note:} The answer correctly identifies Marx's critique of bureaucracies as prioritizing capital over worker welfare, leading to alienation, and connects this to the potential negative impact of similar structures on contemporary business practices, emphasizing the relevance of his ideas in today's organizational contexts.
\end{enumerate}

\vfill
\noindent\textit{End of Answer Key}
\end{document}